% META: {"id":"1SPE-PRODSCAL-CO-025","chapitre":"1SPE-PRODUIT-SCALAIRE","type_objet":"corrige","exercice_id":"1SPE-PRODSCAL-EX-025","status":"approved","origin":"generated","status_history":["generated","audited","accepted","approved"],"release_acceptance":"RELEASE_OWNER_BATCH_ACCEPTANCE","acceptance_closure_digest":"sha256:9b3ccf9a81c5520fbb7e03b7b2d3e7bf2057a02834b6d908bf3be5f49f3b553f","acceptance_receipt_digest":"sha256:4fad86f0282e0fa4e0419cca1ad6379ae7af5a280c04a4a90880f90a1c044242"}
% BEGIN-VERIFY
% from sympy import *
% # A(1,3), B(4,1), C(-1,0)
% ABx, ABy = 3, -2
% ACx, ACy = -2, -3
% ps = 3*(-2) + (-2)*(-3)
% assert ps == 0
% END-VERIFY
\begin{corrige}{1SPE-PRODSCAL-EX-025}
\begin{enumerate}
\item $\overrightarrow{AB}\begin{pmatrix} 3 \\ -2 \end{pmatrix}$ et $\overrightarrow{AC}\begin{pmatrix} -2 \\ -3 \end{pmatrix}$.

$\overrightarrow{AB} \cdot \overrightarrow{AC} = 3 \times (-2) + (-2) \times (-3) = -6 + 6 = 0$.

Le triangle $ABC$ est rectangle en $A$.

\item $AB = \sqrt{9 + 4} = \sqrt{13}$, $AC = \sqrt{4 + 9} = \sqrt{13}$, $BC = \sqrt{25 + 1} = \sqrt{26}$.
\end{enumerate}
\end{corrige}
